\chapter{ТЕОРЕТИЧЕСКИЕ ОСНОВЫ И СУЩЕСТВУЮЩИЕ ПОДХОДЫ}

\section{Системная динамика и проверка моделей}

\subsection{Уровни, потоки и обратные связи}

Системная динамика исследует изменение состояния системы во времени через накопители (stocks), скорости их изменения (flows) и обратные связи. Основы подхода заложены работами Форрестера \cite{forrester1961industrial}, а современное изложение применения к сложным системам дано Стерманом \cite{sterman2000business}. Для набора уровней $S_1,\ldots,S_n$ модель записывается в виде

\begin{equation}
    \frac{dS_i}{dt}=F_i(S_1,\ldots,S_n;P_1,\ldots,P_m;t),\qquad i=1,\ldots,n,
    \label{eq:sd-general}
\end{equation}

где $P_j$ --- параметры, а $F_i$ определяется суммой входящих и исходящих потоков:

\begin{equation}
    F_i=\sum_{r\in I_i}r-\sum_{r\in O_i}r.
    \label{eq:flow-balance}
\end{equation}

Здесь $I_i$ и $O_i$ --- множества входящих и исходящих потоков уровня $S_i$. Даже при простой форме уравнения (\ref{eq:sd-general}) содержательные ошибки могут возникать из-за неверного выбора переменных, неправильного знака потока или неправильных единиц параметра.

CLD отображает качественные причинно-следственные связи и их полярность. Знак петли определяется произведением знаков её рёбер: чётное число отрицательных связей соответствует усиливающей петле, нечётное --- балансирующей. SFD дополнительно фиксирует уровни, потоки и внешние источники/стоки. Поэтому CLD не может служить полной заменой уравнений и топологии потоков, а SFD не гарантирует правильность коэффициентов и поведения.

Валидация системно-динамической модели включает не только запуск интегратора. В работах по model validity рассматриваются структурная проверка, соответствие гипотезе, тестирование поведения и сопоставление с данными \cite{barlas1996formal}. Для текстовой генерации это означает необходимость раздельной оценки нескольких свойств: формальной целостности, согласованности структуры, соответствия описанию и ожидаемого поведения.

\subsection{Размерность и временная сетка}

Пусть каждой величине сопоставлен вектор целочисленных степеней базовых единиц. Для выражений $x$ и $y$ выполняются правила

\begin{align}
    [xy]&=[x]+[y],\\
    [x/y]&=[x]-[y],\\
    [x+y]&=[x]=[y].
\end{align}

Единица потока, соединённого с уровнем $S$, должна быть $[S]/[t]$. Это ограничение позволяет автоматически обнаруживать, например, коэффициент с единицами $minute^{-1}$ в модели, где время задано в годах. Скрытая конвертация между днями, месяцами и годами запрещена: отсутствие conversion registry должно приводить к диагностике, а не к подгонке значения.

Численная симуляция также зависит от временной сетки. При методе Эйлера

\begin{equation}
    S_{k+1}=S_k+hF(S_k,t_k),\qquad t_{k+1}=t_k+h,
    \label{eq:euler}
\end{equation}

где $h$ --- шаг интегрирования. Запрошенный метод, шаг и интервал сохранения результатов должны совпадать с фактически выполненными. Замена Euler на адаптивный метод может изменить близкую к нулю траекторию или временную сетку, поэтому численную совместимость нельзя маскировать более сильным solver.

\section{Генерация структурированных моделей с помощью LLM}

\subsection{Возможности и ограничения}

Трансформер и последующие LLM создали технологическую основу для извлечения структурированных данных из текста \cite{vaswani2017attention}. Few-shot prompting добавляет примеры ожидаемого формата \cite{brown2020gpt3}, метод retrieval-augmented generation (RAG) дополняет контекст внешними источниками \cite{lewis2020retrieval}. Однако ни один из этих подходов не гарантирует корректность формальной модели. Ограничения JSON Schema уменьшают число синтаксических ошибок, но не проверяют смысл переменных или физическую размерность.

В текстовой задаче можно выделить следующие классы ошибок:

\begin{itemize}
    \item неправильный выбор уровней и направления причинно-следственных связей;
    \item нарушение типов и схемы ответа;
    \item потеря или переименование ссылок между этапами;
    \item неправильные коэффициенты и единицы измерения;
    \item добавление незапрошенных механизмов и параметров;
    \item неполное покрытие горизонта, нечисленные значения и вырождение поведения;
    \item нестабильность при повторном запуске.
\end{itemize}

Генерация Python-кода напрямую особенно чувствительна к контексту и несовместимости версий библиотек. Более надёжный путь --- получить ограниченное формальное представление и использовать доверенный генератор кода. Систематические обзоры LLM для code generation описывают как потенциал, так и потребность в проверке сгенерированных артефактов \cite{jiang2024survey}.

\subsection{Связанные работы}

ChatPySD демонстрирует встраивание и симуляцию моделей системной динамики в диалоговом интерфейсе ChatGPT-4 \cite{hu2025chatpysd}. Основное внимание уделено работе с существующими моделями и взаимодействию через PySD. Это отличается от задачи построения новой модели по текстовой гипотезе.

Liu и Keith исследуют генерацию CLD из динамических гипотез \cite{liu2025leveraging}. В их работе применяются curated prompting techniques и сравнение четырёх сочетаний методов с экспертными CLD. Результаты подтверждают возможность получения качественных диаграмм для простых структур, но не решают задачу размерно согласованной исполняемой модели с потоками, начальными значениями и ограниченным repair.

Akhavan и Jalali обсуждают возможности и ограничения генеративного ИИ в simulation modeling \cite{akhavan2024generative}, а обзор Muthiah и др. отображает применение ИИ на разных стадиях системно-динамического моделирования \cite{muthiah2025mapping}. Эти работы показывают, что простое использование LLM не является самостоятельной новизной: научный интерес представляет проверяемый метод, оценка его ограничений и сравнение с альтернативами.

\begin{table}[H]
    \centering
    \caption{Сопоставление направлений исследования}
    \label{tab:related-work}
    \small
    \begin{tabularx}{\textwidth}{L{3cm}L{4.1cm}L{3.3cm}X}
        \toprule
        Работа & Основной предмет & Оценка & Место настоящей работы \\
        \midrule
        ChatPySD & Диалог и симуляция существующих моделей & Демонстрация интеграции & Генерация новой модели и проверка её свойств \\
        Liu и Keith & Текстовая гипотеза $\rightarrow$ CLD & Сравнение prompting techniques с экспертами & Дополнение CLD уравнениями, units, потоками и поведением \\
        Akhavan и Jalali & Возможности и ограничения генеративного ИИ & Аналитический обзор & Конкретный механизм ограничения генерации \\
        Muthiah и др. & Стадии применения ИИ в системной динамике & Обзор направлений & Типизированный план, repair и interoperability \\
        Настоящая работа & Текст $\rightarrow$ исполняемая модель & Controlled ablation и paired tests & Не утверждается превосходство без парного baseline \\
        \bottomrule
    \end{tabularx}
\end{table}

\section{Постановка задачи и требования}

Пусть дано текстовое описание $D$. Необходимо построить документ $M$, содержащий переменные, выражения, время, единицы и топологию, такой что $M$ может быть проверен и симулирован. LLM рассматривается как источник гипотезы о структуре, а не как оракул, чей ответ принимается без проверки.

К методу предъявляются следующие требования:

\begin{enumerate}
    \item \textbf{Единство представления.} Уравнения, графы, экспорт и симуляция должны порождаться из одного канонического документа.
    \item \textbf{Ограниченность генерации.} Типы переменных и выражений ограничены; произвольный пользовательский код не выполняется.
    \item \textbf{Сохранение заданного смысла.} Нулевые начальные значения и явные коэффициенты не заменяются эвристическими значениями.
    \item \textbf{Проверяемость.} Ссылки, units, topology, начальные значения и auxiliary-циклы проверяются до возврата результата.
    \item \textbf{Ограниченное исправление.} Число попыток repair ограничено, а предыдущие уравнения не перезаписываются без причины.
    \item \textbf{Прозрачность вычислений.} Метод интегрирования, временная сетка, единицы и ограничения поддержки фиксируются фактически.
    \item \textbf{Воспроизводимость оценки.} Описания, конфигурация, промежуточные результаты и ошибки сохраняются вместе с идентификаторами версии pipeline.
\end{enumerate}

Таким образом, объектом дальнейшего проектирования является не отдельный prompt, а система ограничений и проверок, внутри которой LLM выполняет ограниченную интерпретирующую роль.
